% Variant: runtimes, engine work, cross-platform infrastructure, open source.
\documentclass{resume}
\resumetitle{Systems Engineer -- Runtimes, Engines and Cross-Platform Tooling}
\resumekeywords{Rust, Lua, Luau, LuaJIT, C++, scripting runtimes, game engine, OpenMW, cross-platform, CI/CD, open source, Linux}
\begin{document}
\makeheader

\section{Summary}
\summary{Systems programmer working in Rust, C++ and Lua/Luau. I build scripting runtimes and engine tooling,
and the build infrastructure that ships them on Linux, Windows, macOS and ARM64. Contributor to the OpenMW
engine; author of the Jess runtime and the l3i Luau binder. I find the constraint holding a project
back and build the smallest system that removes it.}

\section{Technical Capabilities}
\skill{Languages}{Rust, C++, Lua / LuaJIT / Luau, Bash}
\skill{Runtimes and engines}{binding and embedding script VMs (hand-written C API bindings, FFI), JIT hooks,
language front ends, capability models, package formats, asset and plugin pipelines}
\skill{Cross-platform delivery}{Cargo, CMake, clang and lld with cross-language LTO, glibc baselines, ARM64
cross-compilation, GitHub Actions, GitLab CI, Sigstore Cosign}

\section{Selected Engineering Work}

\entry{Jess}{native application kernel for Luau, in Rust}{Independent \textperiodcentered{} open source \textperiodcentered{} 2026}
\entrymeta{\href{https://github.com/DreamWeave-MP/jess}{github.com/DreamWeave-MP/jess} \textperiodcentered{} \href{https://crates.io/crates/jess}{crates.io/crates/jess} \textperiodcentered{} MIT or Apache-2.0 \textperiodcentered{} v0.1.0 (21 Oct 2026) to v0.3.0 (5 Nov 2026)}
\begin{bullets}
  \item Drew one boundary and held it: Rust owns operating-system and hardware contact, resource lifetimes and trust boundaries; Luau owns application semantics and policy.
  \item Shipped the whole toolchain as one binary: runtime, test runner, type checker, package tool and style checker. The first release carried signed archives for Linux x64, macOS ARM64 and x64, Windows x64 and PortMaster ARM64.
  \item Designed an immutable package archive pinned by digest in a lockfile; a workspace packs into one file or a standalone executable. Added authenticated ranged file reads and crash-durable file creation in 0.3.0.
\end{bullets}

\entry{l3i}{Luau runtime binder, in Rust and C++}{Independent \textperiodcentered{} open source \textperiodcentered{} 2026}
\entrymeta{\href{https://github.com/DreamWeave-MP/l3i}{github.com/DreamWeave-MP/l3i} \textperiodcentered{} \href{https://crates.io/crates/l3i}{crates.io/crates/l3i} \textperiodcentered{} 1.0.0 (29 Sep 2026) to 2.7.0 (5 Nov 2026)}
\begin{bullets}
  \item Binds upstream Luau directly instead of through a general-purpose crate: a pinned submodule built by the build script, with the whole C API declared by hand. Began as my C++ Luau binder for OpenMW.
  \item Plans every VM before it exists (userdata tags, compiler type slots, atoms) and renders the matching type definitions from the same plan.
  \item Lowers hot script calls to native code through JIT hooks written in Rust; adds a dialect front end (comprehensions, slices, destructuring) whose tooling reports in original source coordinates.
\end{bullets}

\entry{OpenMW}{open-source reimplementation of the Morrowind engine, C++}{Open source \textperiodcentered{} merged upstream}
\entrymeta{\href{https://gitlab.com/OpenMW/openmw}{gitlab.com/OpenMW/openmw} \textperiodcentered{} described on the project's blog as a core developer (\href{https://openmw.org/2024/from-bsp-to-esp-how-s3ctor-abused-quake-editors-to-redefine-the-morrowind-modding-experience/}{Nov 2024})}
\begin{bullets}
  \item Merged: object movement in the OpenMW-CS editor (\href{https://gitlab.com/OpenMW/openmw/-/merge_requests/3675}{!3675}); unified instance drop modes with an in-editor tooltip (\href{https://gitlab.com/OpenMW/openmw/-/merge_requests/4350}{!4350}); fixes so cell water height is saved whenever the user changed it and journal records are serialized first when saving plugins.
  \item Prototype, not merged: a networking layer and headless server on the \href{https://github.com/DreamWeave-MP/DreamWeave}{DreamWeave fork}, toward upstream multiplayer.
\end{bullets}

\entry{Morrobroom and vfstool}{engine tooling, in Rust}{Independent \textperiodcentered{} open source}
\entrymeta{\href{https://github.com/magicaldave/Morrobroom}{github.com/magicaldave/Morrobroom} \textperiodcentered{} \href{https://github.com/DreamWeave-MP/vfstool}{github.com/DreamWeave-MP/vfstool} \textperiodcentered{} \href{https://crates.io/crates/vfstool}{crates.io/crates/vfstool}}
\begin{bullets}
  \item Morrobroom compiles TrenchBroom (Quake) maps into Morrowind plugins and meshes, turning brush half-spaces into game geometry (featured on the OpenMW blog). vfstool reimplements OpenMW's virtual file system as a library and CLI for Linux, Windows and macOS, on crates.io and the AUR.
\end{bullets}

\entry{StroggForge}{shared build and release infrastructure}{Independent \textperiodcentered{} open source \textperiodcentered{} 2026}
\entrymeta{\href{https://github.com/DreamWeave-MP/StroggForge}{github.com/DreamWeave-MP/StroggForge} \textperiodcentered{} GitHub Actions, Bash}
\begin{bullets}
  \item Reusable workflows and a composite release action shared by DreamWeave's Rust projects: build, portability check, keyless Cosign signing, VirusTotal scan, publish. Linux releases hold a glibc 2.34 baseline; ARM64 builds cross-compile.
\end{bullets}

% Paid employment. Shared by every variant.
\section{Professional Experience}

\entry{Conduent Business Services}{Technical Support Advisor, Tier 2}{Mar 2017 -- Jan 2019}
\begin{bullets}
  \item Tier 2 support for hardware and software across multiple platforms. Owned complex cases from first report to resolution and mentored new team members.
\end{bullets}

\entry{United States Cellular Corporation}{Frontline Customer Service Representative}{2019}
\entrymeta{First-line triage of billing and technical issues, escalating the problems that needed it.}

% Training. Shared by every variant. Neither entry is a degree, and the
% wording must not suggest one.
\section{Education and Training}

\entry{Holberton Tulsa}{AR/VR Engineering and Game Design (20-month program)}{2022 -- 2023}

\entry{Tulsa Tech Riverside}{Computer Repair and Network Technology (A+ certification course)}{2015 -- 2016}


\makefooter
\end{document}
